\chapter{辛几何}

\section{近复结构}

复数可以如下表示，
体现实线性映射、酉群$\text{U}(1)\simeq \text{SO}(2)$：

$$
z = x + y \text i 
= |z| \cdot \text e^{\theta \text i} = \begin{bmatrix}
  x & -y \\ y & x 
\end{bmatrix}
= \sqrt{x^2 + y^2} \begin{bmatrix}
  \cos\theta & -\sin\theta \\ \sin\theta & \cos\theta
\end{bmatrix}
$$

若把$\mathbb C$视为$2 \times 2$实矩阵空间$\mathbb R_2^2 \in\textbf{Vct}_\mathbb R$中的一个二维线性子空间，
取基：

$$
I = \begin{bmatrix}
  1 & 0 \\ 0 & 1
\end{bmatrix}, \quad
J = \begin{bmatrix}
  0 & -1 \\ 1 & 0
\end{bmatrix}, \quad
$$

使得$z = x I + y J$，
且有$J^2 = J \circ J = -I$，
相当于$\text i^2 = -1$。

现在延伸到$n$-维复矩阵空间：

$$
\big( \mathbb C_n^n \in\textbf{Vct}_\mathbb C \big)
\simeq
\big( \mathbb R_{2n}^{2n} \in\textbf{Vct}_\mathbb R \big)
$$

一个复矩阵类似表示为：

$$
\begin{bmatrix}
  z_1^1 & \cdots & z_1^n \\
  \vdots & \ddots & \vdots \\
  z_n^1 & \cdots & z_n^n \\
\end{bmatrix}
=
\begin{bmatrix}
  x_1^1 & -y_1^1 & \cdots & x_1^n & -y_1^n \\
  y_1^1 & x_1^1 & \cdots & y_1^n & x_1^n \\
  \vdots & \vdots & \ddots & \vdots & \vdots \\
  x_n^1 & -y_n^1 & \cdots & x_n^n & -y_n^n \\
  y_n^1 & x_n^1 & \cdots & y_n^n & x_n^n \\
\end{bmatrix}
$$

为了构造分块矩阵，取它等价的表示形式：

$$
\begin{aligned}
\begin{bmatrix}
  z_1^1 & \cdots & z_1^n \\
  \vdots & \ddots & \vdots \\
  z_n^1 & \cdots & z_n^n \\
\end{bmatrix}
&=
\left[
\begin{array}{ccc|ccc}
  x_1^1 & \cdots & x_1^n & -y_1^1 & \cdots & -y_1^n \\
  \vdots & \ddots & \vdots & \vdots & \ddots & \vdots \\
  x_n^1 & \cdots & x_n^n & -y_n^1 & \cdots & -y_n^n \\
  \hline
  y_1^1 & \cdots & y_1^n & x_1^1 & \cdots & x_1^n \\
  \vdots & \ddots & \vdots & \vdots & \ddots & \vdots \\
  y_n^1 & \cdots & y_n^n & x_n^1 & \cdots & x_n^n \\
\end{array}
\right]
\\
&=
\left[
\begin{array}{c|c}
  X & -Y \\
  \hline
  Y & X
\end{array}
\right]
= X I_n + Y J_n
\end{aligned}
$$


类似前述有基：

$$
I = 
\left[
\begin{array}{c|c}
  I_n & 0 \\
  \hline
  0 & I_n
\end{array}
\right]
, \quad
J = 
\left[
\begin{array}{c|c}
  0 & -I_n \\
  \hline
  I_n & 0
\end{array}
\right]
$$

$I$是$\mathbb R_{2n}^{2n}$的代数单位元，
$J$则满足：

$$
J^2 = J \circ J = 
\left[
\begin{array}{c|c}
  0 & -I_n \\
  \hline
  I_n & 0
\end{array}
\right]
\circ
\left[
\begin{array}{c|c}
  0 & -I_n \\
  \hline
  I_n & 0
\end{array}
\right]
=
\left[
\begin{array}{c|c}
  -I_n & 0 \\
  \hline
  0 & -I_n
\end{array}
\right]
= - I.
$$

复数单位$\text i$和基矩阵$J$具有相同的性质：

\begin{framed}
  
实向量空间$V$上的线性变换$J \in \text{End} \ V$若满足

$$
J^2 = -\text{id}_V
$$

称$J$是一个\textbf{近复结构}(almost complex structure)。

\end{framed}

在实矩阵表示下，
实转置对应着复共轭转置，
共轭携带的信息不变。
在不混淆的情况下，
把$J$写为转置/共轭的形式：

$$
J = \left[
\begin{array}{c|c}
0 & I_n\\
\hline
-I_n & 0
\end{array}
\right].
$$

它就是辛矩阵。

\section{辛结构}

作用在$k$-线性空间上，
基于左右元的区分，
一对对偶算子$f$作用在左元，
$f^*$作用在右元．
若$\forall x \in X, \beta \in Y^*:\ \langle x,f^* \beta \rangle = \langle f x, \beta \rangle $，
则互为伴随算子．
这里$\langle -,- \rangle$是$X \times X^*$和$Y \times Y^*$上的自然配对。
考虑$X$上的线性自映射，
自伴算子$A = A^\dagger$，
若引入内积则自然变换可以用内积描述为：

$$
% https://tikzcd.yichuanshen.de/#N4Igdg9gJgpgziAXAbVABwnAlgFyxMJZABgBpiBdUkANwEMAbAVxiRAA0QBfU9TXfIRRkATFVqMWbdgD0AVN14gM2PASIjy4+s1aIOivqsEbSY6jqn7ZCnkYHrhpAIzbJekAB1PDOmADmDDAABAC0pACC8mHB3gBOfoGsdsr8akLImq4W7mzevgFBwRHhobGeCYXJSioOGQAsWjm6bAAehqnGjsgAbE0SLfoR7Sm16UR95gNWIACeHWMmKI1Tlh4RwfOjaUvIjdnTHvmJRa2Rm+WVSQs73X0Ha3k+JyEbZ7OXL9ziMFD+8ERQAAzOIQAC2SDIIBwECQmkObAiIGovgARjAGAAFW5CEBBIE4Dog8FIADM1BhSGczRmUVsSmJEMQjWhsMQxBSjKQLMpiGcnNBTIArBS2SIBSTECLWWSJUy+jLEAB2Cl0LAMNhguhoOCUuVIAAcoqQAE5VerNdrdbD9XyobyFTg1Rr9FqdXqGYKqfa2Wboc7Le6bZ7Jc5qYqVf6La6rR7gV6+eHeUaoy6QG7rUSE84fVTwwwsGAPFAIExUUFkSAABYwOhQNiQIuVp3R8AEZIULhAA
\begin{tikzcd}
X \arrow[rr, "A"]                             &  & X                                          &  & x \arrow[rr, maps to]                                                                         &  & Ax                                                            \\
{\langle -,A^* - \rangle} \arrow[u] \arrow[d] &  & {\langle A-,- \rangle} \arrow[u] \arrow[d] &  & {\langle x,A y \rangle} \arrow[u, maps to] \arrow[d, maps to] \arrow[rr, no head, Rightarrow] &  & {\langle A x,y \rangle} \arrow[u, maps to] \arrow[d, maps to] \\
X^*                                           &  & X^* \arrow[ll, "A^*"]                      &  & A y                                                                                           &  & y \arrow[ll, maps to]                                        
\end{tikzcd}
$$

复线性空间的内积是Hermite内积，
它的实部是实内积，
虚部是辛内积。

线性空间$V$上的双线性形式称为斜对称形式，
若满足：

$$
\omega (u,v) + \omega(v,u) = 0
$$

$2n$-维实线性空间$V$配置非退化斜对称形式$\omega$，
构成辛空间$(V,\omega)$。

若$J$是一个满足
$J^2 = -I$ 的线性变换，则称其为一个\textbf{复结构}。在这种情况下，
可由$J$定义一个由

$$
\omega_J(u,v) := \langle Ju,v\rangle
$$

给出的辛形式。特别地，若$V$配备一个内积$\langle\cdot,\cdot\rangle$，
则一个矩阵$A\in \mathrm{GL}(2n,\mathbb R)$称为\textbf{辛矩阵}，若它满足

$$
A^T J A = J,
$$

在\cite{ChrissGinzburg1997}中如此定义： 

\begin{framed}
    
A symplectic structure on \( X \) is a non-degenerate regular (i.e., \( C^\infty \), resp. holomorphic, algebraic) 2-form \( \omega \) such that \( d\omega = 0 \).

\end{framed}

Let \( X = \mathbb{C}^{2n} \) with coordinates \( q_1, \ldots, q_n, p_1, \ldots, p_n \). Then

\[
\omega = dp_1 \wedge dq_1 + \cdots + dp_n \wedge dq_n
\]

is a symplectic structure.